\section{Flexible Learning Week: 16th - 22nd February 2026}

I spent majority of last week and Monday of this week working on TGP and preparing for my group presentation so my priorities were shifted a little bit. I also had a cold/flu which left me bed ridden for the majority of reading week so I only was able to do work on SHP towards the end of the week.

To-do list for this week:

\begin{itemize}
    \item Have another look at the cylindrical power spectrum graphs from week 4
    \item Look at the GPR notebook (3)
    \item Look at the eigenvectors plot from the paper Alkistis sent. See Figure 27
    \item Have a read of the paper by Zhaoting that uses PCA and the cylindrical power spectrum
    \item Start my report. Create the structure and think of what I want to submit for my 5 pages due next Friday
\end{itemize}

\subsection{Saturday 21st February}

I had a look at the GPR notebook (3)

\begin{itemize}
    \item I ran all the cells in the notebook and noticed there was an error with the "optimize.fmin" function. I tried to figure out what was the source of the error but wasn't sure in the end.
    \item So, I decided to use a different function to find the best fitted parameters which I was more familiar with since I used it in NumRec which was "optimize.minimize" which is also a scipy function.
    \item I tried 3 different optimisations with this function: (1) the default method (BFGS) with no bounds, (2) the default method (BFGS) with bounds, and (3) Nelder-Mead
    \begin{itemize}
        \item When I tried (1) the default method (BFGS) with no bounds, the result of the optimisation was not a success and this was because the desired error was not necessarily achieved due to precision loss and returned a log-marginial likelihood of -1380.158.
    \end{itemize}
    \begin{itemize}
        \item As a result, I tried (2) which included adding bounds. This resulted in a successful optimisation and returned a log-marginial likelihood of -1538.459. 
        \item I also tried (3), the Nelder-Mead method since that is what Zhaoting's original function used which was also successful in optimisation and returned a log-marginial likelihood of -1540.445.
    \end{itemize}
\item From each of these methods, I plotted the kernals for the HI signal and foreground along with the actual HI signal and foreground. I also plotted the 1D and cylindrical power spectrums. Some of the plots, specifically the power spectrum ones, don't look right to me. I am unsure if they are meant to look like that or if I have done something wrong with my code since I did not have any plots to compare to from Zhaoting's code since there was an error.
\item I plotted the resulting graphs from all 3 methods. I also decided to plot the graphs produced from the initial values which are put into the minimize function, just to see what that produced. The graphs from the initial values and all three methods can be seen.
\end{itemize}

I also had another look at the cylindrical power spectrum graphs from week 4

\begin{itemize}
    \item I still don't really understand. I tweaked my code a little bit and since my graph outputs are connected to my overleaf, my plots which are seen in figures 27 - 29 have changed and I don't know why. They will probably continuously change as I edit and run the code. I should change that so I does not continuously change and update like that and save the newly edited code to a new png file.
\end{itemize}

\subsection{Sunday 22nd February}

\begin{itemize}
    \item I had another look at the GPR notebook (3) and the PCA notebook (2). I decided to plot the residuals for both methods so I can better compare the before after after of the foreground cleaning. 
    \item I also noticed that for method (2) for GPR which was 
\end{itemize}